\section{Results}

\subsection{Distance-Dependent Decoding}
If numerosity is represented as a continuous magnitude, conditions that
are numerically farther apart should produce more distinct neural
patterns and thus higher decoding accuracy. We tested this by grouping
condition-pair decoding accuracies by the absolute difference between
the two conditions' target numerosities and computing the group mean at
each time window.
Figure~\ref{fig:temporal-distance}a shows that pairs separated by
greater numerical distance tend to be decoded more accurately, with
the clearest graded ordering by distance in early (80-200 ms) and mid-latency
(340-500 ms) windows. Across all six numerosities, this graded pattern is consistent with a
distance-dependent neural code.

\paragraph{The ``No-1'' control.}
To assess whether the distance gradient is driven by the inclusion of
numerosity~1, we repeated the analysis excluding all pairs
containing ``1'' (Figure~\ref{fig:temporal-distance}b). Sustained
separation by numerical distance largely disappears, with the clearest
residual distance-related gradient limited to approximately 80-120 ms. This suggests that
numerosity~1 disproportionately drives discriminability.

\begin{figure}[!tbp]
  \centering
  \begin{minipage}[t]{\columnwidth}
    \centering
    \includegraphics[width=\columnwidth]{../results/analysis-lda/temporal_change24_all-100-700ms/figures/temporal_distance_timecourse.png}
    \vspace{-2mm}
    \subcaption{All pairs (including numerosity 1)}
  \end{minipage}

  \vspace{2mm}

  \begin{minipage}[t]{\columnwidth}
    \centering
    \includegraphics[width=\columnwidth]{../results/analysis-lda/temporal_change24_all-100-700ms/figures-no1/temporal_distance_timecourse.png}
    \vspace{-2mm}
    \subcaption{Excluding pairs containing numerosity 1}
  \end{minipage}
  \caption{Decoding accuracy by numerical distance ($-$100 to
  700\,ms). Data are from the 24-condition prime-target analysis. Mean
  decoding accuracy (balanced accuracy, $N = 24$) grouped by
  target numerical distance. (a)~With all pairs, decoding tends to increase with
  numerical distance, with the clearest distance-ordered separation in early
  (80-200\,ms) and mid-latency (340-500\,ms) windows. (b)~Excluding
  pairs containing numerosity~1 reduces the graded ordering, with the
  clearest residual gradient limited to approximately 80-120\,ms.}
  \label{fig:temporal-distance}
\end{figure}


\subsection{Whole-Epoch Analysis}
To compare candidate representational structures, we computed
whole-epoch target RDMs and compared them to theoretical models
(Table~\ref{tab:static-model-comparison},
Figure~\ref{fig:static-model-rdms}).

\begin{table}[!t]
  \centering
  \caption{Whole-epoch model correlations (Spearman $\rho$) for the
  target-only analysis. Target-only (1-6) includes all numerosities.
  Target-only no-1 (2-6) excludes numerosity~1 from both prime and
  target during classifier training and evaluation. $p$-values are Holm-corrected across models.
  The boundary contrast tests for a difference in correlation between
  Target PI(x-4) and Target PI(x-3) (paired $t$-test on Fisher-$z$
  differences, two-sided, Cohen's $d_z$). The noise ceiling is a
  descriptive leave-one-subject-out lower-bound estimate of the
  model-brain correlation attainable for structure shared across
  subjects.}
  \label{tab:static-model-comparison}
  \small
  \begin{tabular}{@{}lcccc@{}}
    \toprule
    & \multicolumn{2}{c}{Target-only (1-6)} & \multicolumn{2}{c}{Target-only no-1 (2-6)} \\
    \cmidrule(lr){2-3} \cmidrule(lr){4-5}
    Model & $r$ & $p$ & $r$ & $p$ \\
    \midrule
    Target PI(x-4) & 0.289 & $<0.001$ & 0.269 & 0.003 \\
    Target PI(x-3) & 0.207 & $<0.001$ & 0.076 & ns \\
    ANS            & 0.313 & $<0.001$ & 0.074 & ns \\
    Pixel          & 0.056 & ns        & 0.069 & ns \\
    RT             & 0.162 & 0.005      & 0.268 & 0.007 \\
    Noise ceiling  & 0.417 & \multicolumn{1}{c}{} & 0.332 & \multicolumn{1}{c}{} \\
    \midrule
    \multicolumn{1}{@{}l}{Boundary} & $d$ & $p$ & $d$ & $p$ \\
    \midrule
    PI(x-4) vs PI(x-3) & 0.53 & 0.017 & 0.79 & $<0.001$ \\
    \bottomrule
  \end{tabular}
\end{table}


When all six target numerosities (1-6) were included, ANS, Target
PI(1-4), and Target PI(1-3) were all significant
(Table~\ref{tab:static-model-comparison}),
with ANS showing the highest correlation. Target PI(1-4) outperformed
Target PI(1-3) ($d = 0.53$, $p = 0.017$). The Pixel control did not
reach significance, suggesting that low-level image similarity does not
account for the observed neural dissimilarity structure. RT was significant ($r = 0.162$, $p = 0.005$),
indicating that neural dissimilarity structure covaries with RT
differences.

After excluding numerosity~1, Target PI(2-4) remained significant ($r = 0.269$,
$p_{\text{Holm}} = 0.003$), while Target PI(2-3) ($r = 0.076$) and ANS
($r = 0.074$) did not reach significance. The boundary contrast strengthened to $d = 0.79$ ($p < 0.001$).
The Pixel control remained non-significant, whereas RT remained
significant ($r = 0.268$, $p_{\text{Holm}} = 0.007$).
Among the numerosity models, only Target PI(2-4) remained significant
when numerosity~1 was removed.

\paragraph{Multidimensional scaling.} MDS of the whole-epoch brain RDM
(Figure~\ref{fig:static-change6-mds}) visualizes this structure.
Numerosity~1 is isolated, PI-range
numerosities (2-4) group together, and ANS-range numerosities (5-6)
occupy a separate region, consistent with a categorical boundary at 4-5
and with numerosity~1 as a distinct representational case.

\begin{figure}[!t]
  \centering
  \includegraphics[width=0.85\columnwidth]{../results/analysis-lda/static-change6/static_change6_all/figures-pairwise/mds_category-static_change6_all-pairwise-notitle.png}
  \caption{Multidimensional scaling (MDS) of the whole-epoch
  target-only brain RDM (0-700\,ms, numerosities 1-6). Each point
  represents one target numerosity. Colors indicate PI/ANS category
  under a 4-5 boundary (yellow = 1, blue = PI range 2-4,
  green = ANS range 5-6). Numerosity~1 is isolated from the rest;
  PI-range and ANS-range numerosities occupy separate regions.}
  \label{fig:static-change6-mds}
\end{figure}


\subsection{Time-Resolved Analysis}
%% Queue full-width figures early so LaTeX places them efficiently.
\begin{figure*}[!t]
  \centering
  \includegraphics[width=\textwidth]{../results/analysis-lda/temporal_change24_all-100-700ms/figures/temporal_rdm_snapshots-temporal_change24_all-100-700ms.png}
  \caption{Time-varying neural RDM snapshots for the 24-condition prime-target
  analysis. Each panel shows the mean brain RDM within a 100\,ms bin.
  Conditions are ordered by target numerosity, then prime.}
  \label{fig:temporal-change24-rdm-snapshots}
\end{figure*}


\begin{figure*}[!t]
  \centering
  \begin{minipage}[t]{0.48\textwidth}
    \centering
    \includegraphics[width=\textwidth]{../results/analysis-lda/temporal_change24_all-100-700ms/figures/temporal_model_fits_with_significance-temporal_change24_all-100-700ms-with_RT.png}
    \vspace{-2mm}
    \subcaption{All 24 conditions (with numerosity 1)}
  \end{minipage}
  \hfill
  \begin{minipage}[t]{0.48\textwidth}
    \centering
    \includegraphics[width=\textwidth]{../results/analysis-lda/temporal_change24_all-100-700ms/figures-no1/temporal_model_fits_with_significance-temporal_change24_all-100-700ms-no1-with_RT.png}
    \vspace{-2mm}
    \subcaption{No-1 subset (excluding numerosity 1)}
  \end{minipage}
  \caption{Temporal model fits for the 24-condition prime-target analysis.
  Spearman correlation between neural and model RDMs over time. Colored bars
  indicate significant clusters ($p < 0.05$, cluster permutation).
  (a)~With all pairs, Target PI(1-4) shows two significant clusters
  (50-250\,ms and 270-550\,ms, both $p < 0.001$).
  (b)~After excluding pairs containing numerosity~1, Target PI(2-4)
  retains two significant clusters (50-230\,ms and 250-510\,ms, both
  $p < 0.001$). The gray dotted curve shows the lower-bound noise ceiling.}
  \label{fig:temporal-change24-model-fits}
\end{figure*}

\begin{figure}[!t]
  \centering
  \begin{minipage}[t]{0.48\columnwidth}
    \centering
    \includegraphics[width=\columnwidth,height=0.88\textheight,keepaspectratio]{../results/analysis-lda/temporal_change24_all-100-700ms/figures/temporal_model_rdms_all_compact_col-temporal_change24_all-100-700ms-RT.png}
    \vspace{-2mm}
    \subcaption{All pairs}
  \end{minipage}
  \hfill
  \begin{minipage}[t]{0.48\columnwidth}
    \centering
    \includegraphics[width=\columnwidth,height=0.88\textheight,keepaspectratio]{../results/analysis-lda/temporal_change24_all-100-700ms/figures-no1/temporal_model_rdms_all_compact_col-temporal_change24_all-100-700ms-no1-RT.png}
    \vspace{-2mm}
    \subcaption{No-1 control}
  \end{minipage}
  \caption{Brain and model RDMs for the 24-condition prime-target
  analysis (0-700 ms average). The brain RDM shows mean balanced accuracy
  (\%) across subjects, where each row and column is one prime-target
  condition (e.g., ``42'' = prime 4, target 2). Model RDMs show predicted
  dissimilarity for each theoretical model. (a)~All 24 prime-target
  conditions. (b)~Excluding conditions containing numerosity~1, reducing
  the matrix to 18 conditions.}
  \label{fig:temporal-change24-model-rdms}
\end{figure}

To compare candidate representational structures over time, we
computed time-resolved neural RDMs and compared them to theoretical
models (Figures~\ref{fig:temporal-change24-model-fits}-\ref{fig:temporal-change24-model-rdms};
Appendix Table~\ref{tab:appendix-temporal-model-fit-clusters}).
Figure~\ref{fig:temporal-change24-rdm-snapshots} depicts the evolving
representational structure across all 24 prime-target conditions, with
each panel showing the mean neural RDM in a 100 ms bin.
Figure~\ref{fig:temporal-change24-model-fits} plots the corresponding
time-resolved model-fit trajectories, with bottom bars marking
significant clusters, and
Figure~\ref{fig:temporal-change24-model-rdms} presents the mean brain RDM
(0-700 ms average) alongside the theoretical model RDMs.

When all 24 prime-target conditions were included (with numerosity~1;
Figure~\ref{fig:temporal-change24-model-fits}a), Target PI(1-4) showed
two significant clusters (50-250 ms and 270-550 ms, both $p < 0.001$),
and ANS also showed two significant clusters (30-230 ms and 290-530 ms,
both $p < 0.001$). ANS reached the highest single-window model-brain
correlation in the with-1 condition. Full model-fit cluster windows are
reported in Appendix Table~\ref{tab:appendix-temporal-model-fit-clusters}.

When pairs containing ``1'' were excluded
(Figure~\ref{fig:temporal-change24-model-fits}b), Target PI(2-4)
remained significant in two clusters (50-230 ms and 250-510 ms, both
$p < 0.001$), whereas ANS remained significant in narrower clusters
(50-190 ms, $p < 0.001$, and 410-510 ms, $p = 0.003$). This pattern
indicates that removing numerosity~1 attenuates ANS breadth while
preserving 4-5 boundary structure.

\paragraph{Model-difference tests.}
Table~\ref{tab:model-differences}
shows that the 4-5 boundary captures neural structure better than the
3-4 alternative. Target PI(1-4) significantly outperforms Target PI(1-3)
at 230-330 ms ($p < 0.001$), 490-550 ms ($p = 0.039$), and 570-650 ms
($p = 0.014$). A two-sided Target PI(1-4) vs.\ ANS comparison yields one cluster
at 90-150 ms ($p = 0.050$) in which ANS fits the brain data better
than Target PI (negative $\Delta z$). In the no-1 variant, Target PI(2-4)
exceeds ANS at 170-310 ms ($p < 0.001$) and exceeds Target PI(2-3) at
130-350 ms ($p < 0.001$) and 390-550 ms ($p < 0.001$). Together,
these temporal comparisons converge with the whole-epoch results and
show stronger support for the 4-5 boundary when numerosity~1 is
excluded.

\begin{table}[!tbp]
  \centering
  \caption{Model-difference tests for the 24-condition prime-target analysis
  (cluster permutation, 10,000 permutations, two-sided, $\alpha = 0.05$).
  Best fit indicates which model has the higher brain-model correlation
  in each significant cluster.}
  \label{tab:model-differences}
  \small
  \begin{tabular}{@{}llcr@{}}
    \toprule
    Comparison & Best fit & Range (ms) & $p$ \\
    \midrule
    \multicolumn{4}{l}{\textit{All 24 conditions (with numerosity 1)}} \\
    Target PI(1-4) vs.\ PI(1-3) & PI(1-4) & 230-330 & $<0.001$ \\
    Target PI(1-4) vs.\ PI(1-3) & PI(1-4) & 490-550 & 0.039 \\
    Target PI(1-4) vs.\ PI(1-3) & PI(1-4) & 570-650 & 0.014 \\
    Target PI(1-4) vs.\ ANS     & ANS     & 90-150  & 0.050 \\
    \midrule
    \multicolumn{4}{l}{\textit{No-1 subset (excluding numerosity 1)}} \\
    Target PI(2-4) vs.\ ANS     & PI(2-4) & 170-310 & $<0.001$ \\
    Target PI(2-4) vs.\ PI(2-3) & PI(2-4) & 130-350 & $<0.001$ \\
    Target PI(2-4) vs.\ PI(2-3) & PI(2-4) & 390-550 & $<0.001$ \\
    \bottomrule
  \end{tabular}
\end{table}


\paragraph{Multidimensional scaling.} MDS of the 24-condition neural
RDM (Figure~\ref{fig:temporal-change24-mds}) provides visual
confirmation. Conditions separate by target system membership, with target-1
pairs forming an isolated cluster, PI-range targets (2-4) grouping
centrally, and ANS-range targets (5-6) occupying a distinct region.

\begin{figure}[!t]
  \centering
  \includegraphics[width=\columnwidth]{../results/analysis-lda/temporal_change24_all-100-700ms/figures/temporal_mds_static_0-700ms-change_code-temporal_change24_all-100-700ms.png}
  \caption{MDS projection of the 24-condition neural RDM averaged over
  0-700\,ms. Each point is a prime-target condition code. Colors mark target
  groups under the 4-5 boundary. Yellow marks target~1, blue marks targets
  2-4, and green marks targets 5-6. Target~1 appears separated from the main
  cluster, and targets 2-4 and 5-6 occupy partially distinct regions.}
  \label{fig:temporal-change24-mds}
\end{figure}

